\section{A finite-size criterion for the heat-bath gap}\label{s:knabe}

In this section we bound the heat-bath gap $\lambda(\mu)$ on a torus from below by the smallest heat-bath gap of an $\ell \times \ell$ box. The generator $\sum_x (1 - E_x)$ is a sum of orthogonal projections, two of which commute unless the sites are neighbours, and the box generators play the role of the local Hamiltonians in the finite-size criteria of Knabe \cite{Kna88} and Gosset and Mozgunov \cite{GM16}. We follow the weighted argument of \cite{GM16}, but replace the step that uses a joint eigenvector of the translations by a Cauchy--Schwarz inequality over the placements of the box; in this form the criterion needs no translation invariance of $\mu$, which matters because $\mu_D$ is invariant only under the translations that preserve the checkerboard.

\subsection{Statement}

Let $L, N \ge 3$ and $\Lambda = \Z_L \times \Z_N$, let $(S_x)_{x \in \Lambda}$ be finite sets with $|S_x|\ge2$, and let $\mu$ be a probability measure on $\prod_x S_x$ of the form
\begin{equation}\label{e:MRF}
\mu(s) = \frac{1}{Z} \prod_{x \in \Lambda} \phi_x(s_x) \prod_{\{x, y\} \in E} \phi_{xy}(s_x, s_y), \qquad \phi_x, \phi_{xy} > 0 .
\end{equation}
The frame model $\mu_D$ of \eqref{e:muD} is of this form. Let $E_x$, $\Var_x$ and $\lambda(\mu)$ be as in \eqref{e:lambdaHB}, and set $Q_x = 1 - E_x$ and $\mathcal L = \sum_{x \in \Lambda} Q_x$, a nonnegative operator on $L^2(\mu)$ whose quadratic form is the Dirichlet form. For $\ell \ge 2$ let $R = \{1, \dots, \ell\}^2$, viewed as a subset of $\Lambda$ when $L, N \ge \ell$. For $a \in \Lambda$ and a configuration $\sigma$ on $\Lambda \setminus (R + a)$, let $\gamma(a, \sigma)$ be the spectral gap of $\sum_{x \in R + a} (1 - E_x)$ on $L^2(\mu(\, \cdot \mid \sigma))$, where $E_x$ is the conditional expectation of the conditional measure $\mu(\, \cdot \mid \sigma)$ given the spins of $R + a$ off $x$; it coincides with the restriction of $E_x$. Define
\begin{equation}\label{e:gell}
g(\ell) = \min_{a \in \Lambda} \, \min_{\sigma} \, \gamma(a, \sigma) .
\end{equation}
By the Markov property of \eqref{e:MRF}, $\mu(\, \cdot \mid \sigma)$ depends on $\sigma$ only through the spins on the outer boundary of $R + a$. Finally define
\begin{equation}\label{e:tkappa}
t(\ell) = \frac{36}{5} \, \frac{\ell^2 + 2 \ell + 2}{(\ell + 1)^2 (\ell + 2)^2}, \qquad \kappa_\ell = \frac{25}{36} \, \frac{(\ell + 1)^2 (\ell + 2)^2}{(\ell - 1)(\ell + 3)(\ell^2 + 2 \ell + 2)} .
\end{equation}

\begin{theorem}\label{t:K}
Let $\ell \ge 2$ and $L, N \ge \ell + 1$. Then
\begin{equation}\label{e:K}
\lambda(\mu) \ge \frac{\ell}{\ell - 1} \Bigl( g(\ell) - \frac{1}{\ell} \Bigr) \qquad \text{and} \qquad \lambda(\mu) \ge \kappa_\ell \bigl( g(\ell) - t(\ell) \bigr) .
\end{equation}
\end{theorem}

Only the first bound is used in the proof of Theorem~\ref{t:main}. The second, weighted bound has the threshold $t(\ell) \approx 7.2 / \ell^2$ in place of $1/\ell$; we include it because it is the form in which the criterion is closest to being checkable at moderate $\ell$ (for instance $t(8) = 0.0729$).

\subsection{Three lemmas}

\begin{lemma}\label{l:commute}
If $x \neq y$ are not neighbours, then $E_x E_y = E_y E_x$, and $Q_x Q_y$ is an orthogonal projection; in particular $Q_x Q_y \ge 0$.
\end{lemma}

\begin{proof}
By \eqref{e:MRF}, the conditional law of $s_x$ given the other spins depends only on $s_{N(x)}$, and $y \notin N(x) \cup \{x\}$. Hence
\begin{equation*}
E_x E_y F(s) = \sum_{t \in S_x} \sum_{t' \in S_y} \mu(t \mid s_{N(x)}) \, \mu(t' \mid s_{N(y)}) \, F(s^{x \to t, \, y \to t'}),
\end{equation*}
where $s^{x \to t, \, y \to t'}$ is $s$ with $s_x, s_y$ replaced by $t, t'$. This is symmetric in $x$ and $y$. The product of two commuting orthogonal projections is an orthogonal projection.
\end{proof}

\begin{lemma}\label{l:fibre}
For $a \in \Lambda$ let $A_a = \sum_{x \in R + a} Q_x$. Then $A_a^2 \ge g(\ell) A_a$, and $A_a \ge g(\ell)$ on $(\ker A_a)^\perp$.
\end{lemma}

\begin{proof}
Decompose $L^2(\mu)$ orthogonally according to the values $\sigma$ of the spins on $\Lambda \setminus (R + a)$. For $x \in R + a$, $Q_x$ preserves each summand and acts on the summand of $\sigma$, which is isometric to $L^2(\mu(\, \cdot \mid \sigma))$, as $1 - E_x$. So $A_a$ is the direct sum over $\sigma$ of nonnegative operators whose spectra lie in $\{0\} \cup [\gamma(a, \sigma), \infty) \subset \{0\} \cup [g(\ell), \infty)$. Both claims follow from the spectral theorem.
\end{proof}

For a sequence $u \colon \{1, \dots, \ell\} \to [0, \infty)$, extended by $0$ to $\Z$, let $r(k) = \sum_{j \in \Z} u(j) u(j + k)$, and for $L \ge \ell + 1$ let $r_L(k) = \sum_{j \in \Z_L} \tilde u(j) \tilde u(j + k)$, $k \in \Z_L$, where $\tilde u$ is $u$ on the residues of $1, \dots, \ell$ and $0$ elsewhere.

\begin{lemma}\label{l:autocorr}
Let $u$ be integer valued, symmetric ($u(j) = u(\ell + 1 - j)$) and unimodal, and let $L \ge \ell + 1$. Then $r$ is nonincreasing on $\{0, 1, 2, \dots\}$, and
\begin{equation}\label{e:autocorr}
r_L(0) = r(0), \qquad r_L(\pm 1) = r(1), \qquad r_L(k) \le r(1) \quad \text{for } k \in \Z_L \setminus \{0, \pm 1\} .
\end{equation}
\end{lemma}

\begin{proof}
For $h \ge 1$ let $I_h = \{ j : u(j) \ge h\}$. Since $u$ is symmetric and unimodal, every $I_h$ is an integer interval of the form $\{m, \dots, \ell + 1 - m\}$, whose length is congruent to $\ell$ modulo $2$, and $u = \sum_{h \ge 1} \mathbf 1_{I_h}$ because $u$ is integer valued. Hence $r(k) = \sum_{h, h'} T_{h h'}(k)$ with $T_{h h'}(k) = |I_h \cap (I_{h'} - k)|$. For two integer intervals with the same centre and lengths $p, p'$ of the same parity,
\begin{equation}\label{e:T}
|I_h \cap (I_{h'} - k)| = \min(m_0, \, h_0 - |k|)_+, \qquad m_0 = \min(p, p'), \quad h_0 = (p + p')/2 .
\end{equation}
Each $T_{hh'}$ is nonincreasing in $|k|$, so $r$ is nonincreasing on $\{0, 1, 2, \dots\}$. Note that $h_0 \le \max(p, p') \le \ell$ and $2 h_0 \le m_0 + \ell$.

We now show that $T(k) + T(\ell + 1 - k) \le T(1)$ for $2 \le k \le \ell - 1$, for every $T = T_{hh'}$. Write $x = \ell + 1 - k$; by symmetry we may assume that $k \le x$. If $T(x) = 0$, then $T(k) \le T(2) \le T(1)$. Otherwise $x < h_0$. If $T(k) = m_0$, then $k \le h_0 - m_0 \le (\ell - m_0)/2$, so $x = \ell + 1 - k \ge (\ell + m_0)/2 + 1 \ge h_0 + 1$, which is impossible. Hence $T(k) = h_0 - k$, and
\begin{equation*}
T(k) + T(x) \le 2 h_0 - \ell - 1 \le \min(m_0, \, h_0 - 1) = T(1),
\end{equation*}
where we used $h_0 \le \ell$ and $2 h_0 \le m_0 + \ell$. Summing over $h, h'$ gives $r(k) + r(\ell + 1 - k) \le r(1)$ for $2 \le k \le \ell - 1$.

For $k \in \{0, \dots, L - 1\}$ we have $r_L(k) = r(k) + r(L - k)$, since $j + k$ and $j + k - L$ cannot both lie in $\{1, \dots, \ell\}$; moreover $r(m) = 0$ for $|m| \ge \ell$. This gives $r_L(0) = r(0)$ and $r_L(1) = r_L(L - 1) = r(1)$, because $L - 1 \ge \ell$. Let $2 \le k \le L - 2$. If $k \ge \ell$ or $L - k \ge \ell$, one term vanishes and the other is at most $r(2) \le r(1)$. Otherwise $k \le \ell - 1$ and $\ell + 1 - k \le L - k \le \ell - 1$, so $r(L - k) \le r(\ell + 1 - k)$ by monotonicity, and $r_L(k) \le r(k) + r(\ell + 1 - k) \le r(1)$.
\end{proof}

The flat sequence $u \equiv 1$ and the sequence $u(j) = j(\ell + 1 - j)$ satisfy the hypotheses of Lemma~\ref{l:autocorr}. For the flat sequence $r(0) = \ell$ and $r(1) = \ell - 1$. For $u(j) = j(\ell + 1 - j)$ a direct summation gives
\begin{equation}\label{e:sums}
\sum_{j} u(j) = \frac{\ell(\ell + 1)(\ell + 2)}{6}, \quad r(0) = \frac{\ell(\ell + 1)(\ell + 2)(\ell^2 + 2\ell + 2)}{30}, \quad r(1) = \frac{(\ell - 1)\ell(\ell + 1)(\ell + 2)(\ell + 3)}{30} .
\end{equation}

\subsection{Proof of Theorem~\ref{t:K}}

Fix one of the two sequences $u$, let $w(z) = u(z_1) u(z_2)$ for $z \in R$, $\mathsf s = \sum_j u(j)$, and
\begin{equation*}
B_a = \sum_{z \in R} w(z) \, Q_{z + a}, \qquad a \in \Lambda .
\end{equation*}
For the flat sequence $B_a = A_a$. Since every site lies in exactly $\ell^2$ translates of $R$, we have $\sum_{a} A_a = \ell^2 \mathcal L$ and $\sum_a B_a = \mathsf s^2 \mathcal L$. Expanding the squares and using $Q_x^2 = Q_x$, we have
\begin{equation}\label{e:sumB2}
\mathcal L^2 = \mathcal L + \sum_{x \neq y} Q_x Q_y, \qquad \sum_{a \in \Lambda} B_a^2 = W(0) \mathcal L + \sum_{x \neq y} W(y - x) \, Q_x Q_y,
\end{equation}
where $W(d) = \sum_{v \in \Lambda} \tilde w(v) \tilde w(v + d) = r_L(d_1) r_N(d_2)$, with $\tilde w$ the extension of $w$ by zero. By Lemma~\ref{l:autocorr}, $W(0) = r(0)^2$, $W(d) = r(0) r(1)$ for the four neighbours $d$ of $0$, a number we denote by $W_1$, and $W(d) \le W_1$ for all other $d \neq 0$: if $d_1, d_2 \neq 0$ then $W(d) \le r(1)^2$, and if exactly one coordinate vanishes then the other is not $\pm 1$ and $W(d) \le r(0) r(1)$. By \eqref{e:sumB2},
\begin{equation*}
\mathcal L^2 - \mathcal L - W_1^{-1} \Bigl( \sum_a B_a^2 - W(0) \mathcal L \Bigr) = \sum_{\substack{x \neq y \\ x \not \sim y}} \Bigl( 1 - \frac{W(y - x)}{W_1} \Bigr) Q_x Q_y \ge 0,
\end{equation*}
where we used that the terms of neighbouring pairs cancel and that $Q_x Q_y \ge 0$ for the other pairs by Lemma~\ref{l:commute}. Hence
\begin{equation}\label{e:Lsquare}
\mathcal L^2 \ge \mathcal L + W_1^{-1} \Bigl( \sum_a B_a^2 - W(0) \mathcal L \Bigr).
\end{equation}

\emph{Flat weights.} By Lemma~\ref{l:fibre}, $\sum_a A_a^2 \ge g(\ell) \sum_a A_a = g(\ell) \ell^2 \mathcal L$. Inserting this into \eqref{e:Lsquare} with $W(0) = \ell^2$ and $W_1 = \ell(\ell - 1)$, we obtain
\begin{equation*}
\mathcal L^2 \ge \Bigl( 1 + \frac{\ell^2 (g(\ell) - 1)}{\ell(\ell - 1)} \Bigr) \mathcal L = \frac{\ell}{\ell - 1} \Bigl( g(\ell) - \frac1\ell \Bigr) \mathcal L .
\end{equation*}

\emph{Weighted weights.} Let $\varphi \in L^2(\mu)$, and set $X_a = \langle \varphi, B_a \varphi \rangle$ and $Y_a = \langle \varphi, A_a \varphi \rangle$. If $\langle\varphi,\mathcal L\varphi\rangle=0$, positivity gives $Q_x\varphi=0$ for every $x$, hence $A_a\varphi=B_a\varphi=0$ for every $a$; the required bound
\begin{equation*}
\Bigl\langle\varphi,\sum_aB_a^2\varphi\Bigr\rangle
\ge g(\ell)\frac{\mathsf s^4}{\ell^2}
\langle\varphi,\mathcal L\varphi\rangle
\end{equation*}
then holds with both sides zero. In the remaining calculation assume $\langle\varphi,\mathcal L\varphi\rangle>0$, so $\sum_aY_a=\ell^2\langle\varphi,\mathcal L\varphi\rangle>0$. Let $G_a$ be the orthogonal projection onto $\ker A_a$. Since $Q_x \ge 0$ and $w > 0$ on $R$, $\ker B_a = \bigcap_{x \in R + a} \ker Q_x = \ker A_a$, so $B_a = G_a^\perp B_a G_a^\perp$. If $G_a^\perp \varphi = 0$ then $X_a = Y_a = 0$. Otherwise, by Cauchy--Schwarz and Lemma~\ref{l:fibre},
\begin{equation*}
X_a^2 = \langle G_a^\perp \varphi, B_a \varphi \rangle^2 \le \|G_a^\perp \varphi\|^2 \|B_a \varphi\|^2, \qquad Y_a \ge g(\ell) \|G_a^\perp \varphi\|^2,
\end{equation*}
so $\langle \varphi, B_a^2 \varphi \rangle \ge g(\ell) X_a^2 / Y_a$. By Cauchy--Schwarz over $a$,
\begin{equation*}
\Bigl\langle \varphi, \sum_a B_a^2 \varphi \Bigr\rangle \ge g(\ell) \sum_{a : Y_a > 0} \frac{X_a^2}{Y_a} \ge g(\ell) \frac{(\sum_a X_a)^2}{\sum_a Y_a} = g(\ell) \frac{\mathsf s^4}{\ell^2} \langle \varphi, \mathcal L \varphi \rangle .
\end{equation*}
Inserting this into \eqref{e:Lsquare}, we find for all $\varphi$
\begin{equation*}
\langle \varphi, \mathcal L^2 \varphi \rangle \ge \Bigl( 1 + \frac{g(\ell) \mathsf s^4 / \ell^2 - r(0)^2}{r(0) r(1)} \Bigr) \langle \varphi, \mathcal L \varphi \rangle = \frac{\mathsf s^4}{\ell^2 r(0) r(1)} \Bigl( g(\ell) - \frac{\ell^2 r(0) (r(0) - r(1))}{\mathsf s^4} \Bigr) \langle \varphi, \mathcal L \varphi \rangle .
\end{equation*}
By \eqref{e:sums}, the prefactor is $\kappa_\ell$ and the subtracted term is $t(\ell)$.

\emph{Conclusion.} In both cases $\mathcal L^2 \ge c \mathcal L$ with $c$ the right side of the corresponding bound in \eqref{e:K}. If $c \le 0$ there is nothing to prove. Otherwise every nonzero eigenvalue $\lambda$ of $\mathcal L$ satisfies $\lambda^2 \ge c \lambda$, that is, $\lambda \ge c$. Since $\mu$ has full support, the heat-bath chain is irreducible and $\ker \mathcal L$ consists of the constants, so $\lambda(\mu) \ge c$. \qed

\begin{remark}\label{r:Kpins}
In \eqref{e:gell} the minimum runs over the boundary values that actually occur on the torus. When $L = \ell + 1$, the column of $\Lambda$ outside $R + a$ is adjacent to both vertical sides of the box, so the left and right boundary spins of a row coincide. In \S\ref{s:mixing} we bound $\gamma(a, \sigma)$ below by the gap of the box measure on $\Z^2$ with an arbitrary boundary condition, which covers these cases.
\end{remark}
